\documentclass[11pt]{article}

% Manuscript skeleton only: argument, formal results, and figure plan.
% Full prose, proofs, citations, and production formatting are intentionally deferred.

\usepackage[margin=1in]{geometry}
\usepackage{amsmath,amssymb,amsthm,bm}
\usepackage{booktabs}
\usepackage{enumitem}
\usepackage[hidelinks]{hyperref}
\usepackage{microtype}

\newtheorem{theorem}{Theorem}
\newtheorem{proposition}[theorem]{Proposition}
\newtheorem{corollary}[theorem]{Corollary}
\newtheorem{remark}[theorem]{Remark}

\newcommand{\E}{\mathbb{E}}
\newcommand{\R}{\mathbb{R}}
\newcommand{\tr}{\operatorname{tr}}
\newcommand{\Var}{\operatorname{Var}}
\newcommand{\Cov}{\operatorname{Cov}}
\newcommand{\bhat}{\widehat{\bm\beta}}
\newcommand{\bstar}{\bm\beta^{\star}}
\newcommand{\that}{\widehat{\bm\theta}}
\newcommand{\tstar}{\bm\theta^{\star}}

\title{Prediction Does Not Imply Control:\\
Geometry and Random-Matrix Theory of Self-Generated Distribution Shift}
\author{Author list to be added}
\date{}

\begin{document}
\maketitle

\begin{abstract}
\begin{itemize}[leftmargin=1.5em]
  \item Predictive models of neural responses can have similar held-out
        accuracy yet differ sharply in their ability to synthesize stimuli
        that control the true response.
  \item We formalize this dissociation in a noisy linear teacher--student
        model. Natural generalization is measured in the data-covariance
        geometry, whereas gradient-based accentuation is evaluated along a
        direction chosen by the fitted model itself.
  \item Random-matrix deterministic equivalents quantify generalization,
        accentuation, peer review, and cross-validated ridge selection under
        anisotropic data.
  \item For regression after a fixed linear feature map, prediction depends
        on a smaller set of summary statistics than pixel-space control. We
        construct continuous families of feature maps with identical
        prediction and cross-validation behavior but dramatically different
        accentuation.
  \item Natural-image covariance experiments validate the theory and show
        that moving prediction-equivalent features into low-variance spectral
        directions can leave prediction intact while driving accentuation
        $R^2$ from near one to very negative values.
\end{itemize}
\end{abstract}

\section*{Status and intended scope}

\begin{itemize}
  \item This document is an author-facing manuscript skeleton, not a complete
        draft.
  \item Bullets specify the intended logical content of each paragraph.
  \item Formal statements identify the results that should survive into the
        main text. Proofs and long deterministic-equivalent calculations will
        be placed in the supplement.
  \item Citations, biological details, author list, and journal formatting are
        placeholders.
\end{itemize}

\section{Introduction: the prediction--control paradox}
\label{sec:introduction}

\begin{itemize}
  \item Start from the empirical observation: encoding models with comparable
        held-out prediction can generate qualitatively different preferred
        images and can differ sharply in closed-loop control.
  \item Highlight the evaluation asymmetry: a poor generator can sometimes
        evaluate stimuli produced by a better generator even though its own
        synthesized stimuli fail under the ground-truth response.
  \item State the conceptual reason. Standard generalization evaluates a
        predictor on an externally specified natural distribution. Control
        evaluates it on a distribution selected by the fitted predictor
        itself.
  \item State the paper's thesis: prediction and gradient-based control depend
        on different geometries and, in feature models, on different
        identifiable summary statistics.
  \item Explain the role of random-matrix theory: the geometric gap is exact,
        while RMT predicts how finite samples, response noise, anisotropy,
        ridge shrinkage, and cross-validation determine its magnitude.
  \item Preview the strongest result: a continuous family of feature maps can
        have identical Gaussian training distributions and identical
        prediction-optimal ridge penalties, while its pixel-space control
        changes by orders of magnitude.
\end{itemize}

\begin{figure}[t]
  \centering
  \fbox{\parbox{0.93\linewidth}{
    \textbf{Figure 1 placeholder: empirical paradox and minimal geometry.}
    Left: representative natural-image prediction, successful versus failed
    accentuation, and peer-review asymmetry. Right: natural inputs probe the
    covariance ellipsoid, whereas accentuation moves along the fitted gradient.
  }}
  \caption{Similar prediction does not guarantee similar control because the
  model changes the distribution on which it is evaluated.}
  \label{fig:empirical-paradox}
\end{figure}

\section{Teacher--student setup and evaluation metrics}
\label{sec:setup}

\begin{itemize}
  \item Let centered inputs and noisy responses follow
  \begin{equation}
    \bm x\sim\mathcal N(0,\Sigma),\qquad
    y=\bstar^{\top}\bm x+\sigma\epsilon,
    \qquad \epsilon\sim\mathcal N(0,1).
    \label{eq:teacher}
  \end{equation}
  \item Define the noiseless signal variance
  \begin{equation}
    S=\bstar^{\top}\Sigma\bstar.
  \end{equation}
  \item Distinguish evaluation against the noiseless teacher from evaluation
        against a newly observed noisy response. The main text uses the
        noiseless teacher; test noise can be added as an irreducible term.
  \item For a fixed fitted pixel-space weight $\bhat$, define natural
        generalization and own-gradient accentuation from the same underlying
        teacher.
\end{itemize}

\begin{proposition}[Exact prediction and own-path control metrics]
\label{prop:exact-metrics}
For a fixed fitted weight $\bhat$, natural generalization satisfies
\begin{align}
  E_{\rm gen}
    &=(\bhat-\bstar)^{\top}\Sigma(\bhat-\bstar),\label{eq:egen-exact}\\
  R_{\rm gen}^{2}
    &=1-\frac{E_{\rm gen}}{S}.\label{eq:rgen-exact}
\end{align}
Along the self-generated path $\bm x(t)=\bm x_0+t\bhat$, define
\begin{equation}
  N_{\rm acc}=\bhat^{\top}\bstar,
  \qquad
  D_{\rm acc}=\bhat^{\top}\bhat.
\end{equation}
The true-on-predicted path slope, student-range-normalized accentuation loss,
and true-path coefficient of determination are
\begin{align}
  b_{\rm acc}
    &=\frac{N_{\rm acc}}{D_{\rm acc}},\label{eq:bacc-exact}\\
  \frac{E_{\rm acc}}{S}
    &=\left(1-\frac{N_{\rm acc}}{D_{\rm acc}}\right)^2,\label{eq:eacc-exact}\\
  R_{\rm acc}^{2}
    &=1-\left(1-\frac{D_{\rm acc}}{N_{\rm acc}}\right)^2.\label{eq:racc-exact}
\end{align}
\end{proposition}

\begin{itemize}
  \item Explain the two opposite singularities. $E_{\rm acc}/S$ is unstable
        when the fitted gradient norm $D_{\rm acc}$ approaches zero;
        $R_{\rm acc}^{2}$ is unstable when the true path gain
        $N_{\rm acc}$ approaches zero.
  \item Emphasize that a very negative $R_{\rm acc}^{2}$ need not mean an
        unbounded raw pathwise MSE: when $N_{\rm acc}\to0$ with
        $D_{\rm acc}>0$, $R_{\rm acc}^{2}\to-\infty$ while
        $E_{\rm acc}/S\to1$.
  \item Use $b_{\rm acc}$ for slope and
        $z_{\rm acc}=D_{\rm acc}/N_{\rm acc}-1$ for the $R^2$-controlling
        ratio. Do not call either quantity $R_{\rm det}$.
\end{itemize}

\section{Direct pixel regression: anisotropy, noise, and ridge}
\label{sec:pixel-ridge}

\begin{itemize}
  \item Fit normalized ridge regression using $n$ samples:
  \begin{equation}
    \bhat
    =(\widehat\Sigma+\lambda I)^{-1}
      \left(\widehat\Sigma\bstar+
      \frac{\sigma}{n}X^{\top}\bm\eta\right),
    \qquad \widehat\Sigma=\frac{X^{\top}X}{n}.
    \label{eq:pixel-ridge}
  \end{equation}
  \item Make the ridge convention visible: software minimizing
        $\lVert\bm y-X\bm\beta\rVert^2+alpha\lVert\bm\beta\rVert^2$
        uses $\alpha=n\lambda$.
  \item Introduce the effective deterministic-equivalent penalty $\kappa$:
  \begin{equation}
    \kappa-\lambda
      =\frac{\kappa}{n}\tr\!\left[
       \Sigma(\Sigma+\kappa I)^{-1}\right].
    \label{eq:kappa-fixed-point}
  \end{equation}
  \item Define compact spectral summaries
  \begin{align}
    B_{p,q}(\kappa)
      &=\bstar^{\top}\Sigma^p(\Sigma+\kappa I)^{-q}\bstar,\\
    \mathrm{df}_{p,q}(\kappa)
      &=\tr\!\left[\Sigma^p(\Sigma+\kappa I)^{-q}\right],
  \end{align}
  with $\mathrm{df}_{2}=\mathrm{df}_{2,2}$.
\end{itemize}

\begin{proposition}[Pixel-space deterministic equivalents]
\label{prop:pixel-de}
In the proportional high-dimensional limit, the leading deterministic
equivalent of noiseless generalization error is
\begin{equation}
  E_{\rm gen}^{\rm DE}
  =\frac{n}{n-\mathrm{df}_{2}}
     \left(\kappa^2 B_{1,2}+\sigma^2\right)-\sigma^2.
  \label{eq:pixel-egen-de}
\end{equation}
The leading means of the accentuation numerator and denominator are
\begin{align}
  \mu_N
    &=B_{1,1},\label{eq:pixel-mu-n}\\
  \mu_D
    &=B_{2,2}
      +\frac{\mathrm{df}_{1,2}}{n-\mathrm{df}_{2}}
        \left(\kappa^2B_{1,2}+\sigma^2\right)\label{eq:pixel-mu-d}\\
    &=B_{2,2}
      +\frac{\mathrm{df}_{1,2}}{n}
       \left(E_{\rm gen}^{\rm DE}+\sigma^2\right).
\end{align}
Consequently, the typical-fit leading approximation is
\begin{equation}
  z_{\rm acc}^{\rm DE}
  =\frac{\mu_D}{\mu_N}-1
  =\frac{\sigma^2\mathrm{df}_{1,2}-n\lambda B_{1,2}}
         {(n-\mathrm{df}_{2})B_{1,1}}.
  \label{eq:pixel-z-de}
\end{equation}
\end{proposition}

\begin{itemize}
  \item Interpret the per-PC error decomposition in one paragraph or inset:
        ridge bias suppresses teacher signal below the $\kappa$ scale, while
        finite-sample noise is largest in low-variance modes near that scale.
  \item Explain the central dissociation: $E_{\rm gen}$ weights coefficient
        error by $\Sigma$, whereas $D_{\rm acc}$ measures unweighted pixel
        gradient energy.
  \item Keep the one- and two-resolvent derivations in the supplement.
\end{itemize}

\begin{proposition}[Prediction-optimal ridge need not optimize control]
\label{prop:pixel-cv}
At an interior stationary point of the DE generalization curve,
\begin{equation}
  (n-\mathrm{df}_{2})\kappa B_{2,3}
  =\mathrm{df}_{2,3}\left(\kappa^2B_{1,2}+\sigma^2\right).
  \label{eq:pixel-cv-stationary}
\end{equation}
Substitution into Eq.~\eqref{eq:pixel-z-de} gives
\begin{equation}
  \boxed{
  z_{\rm acc,CV}^{\rm DE}
  =\kappa\frac{\mathrm{df}_{1,2}}{B_{1,1}}
   \left(
     \frac{B_{2,3}}{\mathrm{df}_{2,3}}
     -\frac{B_{1,2}}{\mathrm{df}_{1,2}}
   \right).}
  \label{eq:pixel-z-cv}
\end{equation}
The right-hand side is not generically zero. Thus minimizing natural
prediction risk does not in general calibrate the self-generated path.
\end{proposition}

\begin{itemize}
  \item Separate the population/DE-CV optimum in
        Proposition~\ref{prop:pixel-cv} from random fold-to-fold selection of
        $\alpha$. The latter is a finite-sample fluctuation analyzed in the
        supplement.
  \item Use a dense logarithmic $\alpha$ grid in all experiments so apparent
        transitions are not artifacts of coarse regularization stages.
\end{itemize}

\begin{figure}[t]
  \centering
  \fbox{\parbox{0.93\linewidth}{
    \textbf{Figure 2 placeholder: pixel-ridge mechanism and RMT validation.}
    (a) Data spectrum and teacher power. (b) Mean per-PC coefficient error,
    emphasizing low-variance directions. (c) Noise sweep with DE and Monte
    Carlo for $R_{\rm gen}^{2}$ and $R_{\rm acc}^{2}$. (d) DE-CV and empirical
    CV ridge penalties. Optional peer-review panel.
  }}
  \caption{RMT predicts the prediction--accentuation gap after
  cross-validated ridge selection.}
  \label{fig:pixel-rmt}
\end{figure}

\section{Linear features create two geometries}
\label{sec:feature-geometry}

\begin{itemize}
  \item Map pixels to $p$ fixed linear features and fit ridge in feature
        coordinates:
  \begin{equation}
    \bm z=F\bm x,\qquad
    \widehat f(\bm x)=\that^{\top}F\bm x,
    \qquad \bhat=F^{\top}\that.
    \label{eq:feature-model}
  \end{equation}
  \item Define the prediction summaries
  \begin{equation}
    C=F\Sigma F^{\top},
    \qquad \bm h=F\Sigma\bstar,
  \end{equation}
  and the pixel-control summaries
  \begin{equation}
    G=FF^{\top},
    \qquad \bm q=F\bstar.
  \end{equation}
  \item Explain the distinction: $C$ is the natural feature covariance,
        while $G$ measures the pixel norm of a back-propagated feature-space
        direction. Likewise, $\bm h$ is the feature--teacher covariance on
        natural inputs, while $\bm q$ is the teacher gain along a pixel
        gradient.
\end{itemize}

\begin{proposition}[Prediction-optimal and control-optimal feature weights]
\label{prop:two-centers}
The best population predictor in the feature span and its irreducible error
are
\begin{align}
  \tstar&=C^{\dagger}\bm h,\label{eq:theta-p}\\
  S_{\perp}&=S-\bm h^{\top}C^{\dagger}\bm h.\label{eq:sperp}
\end{align}
The feature coefficient whose back-projected pixel weight is closest to the
teacher in Euclidean norm is
\begin{equation}
  \bm\theta_{\rm E}=G^{\dagger}\bm q.
  \label{eq:theta-e}
\end{equation}
Prediction training centers the estimator at $\tstar$ in the $C$ geometry,
whereas exact pixel-gradient control is centered at $\bm\theta_{\rm E}$ in the
$G$ geometry. These centers need not agree.
\end{proposition}

\begin{proposition}[Feature-space DE has the same structure as pixel-space DE]
\label{prop:feature-de}
Write the feature covariance, prediction-optimal realizable weight, and
unrealizable signal variance as
\begin{align}
  \Sigma_F&:=F\Sigma F^{\top}=C,\\
  \bm w^{\star}&:=\Sigma_F^{\dagger}F\Sigma\bstar=\tstar,\\
  S_{\perp}&=S-(\bm w^{\star})^{\top}\Sigma_F\bm w^{\star}.
  \label{eq:feature-effective-problem}
\end{align}
Conditional on the retained features, feature regression is therefore the
same ridge problem as pixel regression under the substitutions
\begin{equation}
  \boxed{
  \Sigma\ \longmapsto\ \Sigma_F,
  \qquad
  \bstar\ \longmapsto\ \bm w^{\star},
  \qquad
  \sigma^2\ \longmapsto\ \sigma^2+S_{\perp}.}
  \label{eq:pixel-to-feature-substitution}
\end{equation}
Define the feature-space versions of the same spectral summaries,
\begin{align}
  B^{F}_{p,q}(\kappa)
    &=(\bm w^{\star})^{\top}
      \Sigma_F^{p}(\Sigma_F+\kappa I)^{-q}\bm w^{\star},\\
  \mathrm{df}^{F}_{p,q}(\kappa)
    &=\tr\!\left[
      \Sigma_F^{p}(\Sigma_F+\kappa I)^{-q}\right],
  \qquad
  \mathrm{df}^{F}_{2}=\mathrm{df}^{F}_{2,2}.
  \label{eq:feature-B-df}
\end{align}
Here $\kappa$ obeys the same fixed-point equation,
\begin{equation}
  \kappa-\lambda
  =\frac{\kappa}{n}\mathrm{df}^{F}_{1,1}(\kappa).
  \label{eq:feature-kappa-fixed-point}
\end{equation}
The generalization error is consequently
\begin{equation}
  \boxed{
  E_{\rm gen}^{\rm DE}
  \asymp
  \frac{n}{n-\mathrm{df}^{F}_{2}}
  \left[
    \kappa^2B^{F}_{1,2}
    +S_{\perp}+\sigma^2
  \right]-\sigma^2.}
  \label{eq:feature-egen-de}
\end{equation}
This is exactly Eq.~\eqref{eq:pixel-egen-de}, with the substitutions in
Eq.~\eqref{eq:pixel-to-feature-substitution}, except that $S_{\perp}$ remains
part of the noiseless test error because it is omitted teacher signal rather
than measurement noise.

To express pixel-space accentuation, define the shrinkage effective weight
mapped back to input space,
\begin{align}
  \widetilde{\bm\beta}_{\kappa}
  &:=F^{\top}\Sigma_F(\Sigma_F+\kappa I)^{-1}\bm w^{\star}\\
  &=F^{\top}(\Sigma_F+\kappa I)^{-1}F\Sigma\bstar.
  \label{eq:feature-shrunken-pixel-weight}
\end{align}
Then the leading numerator and denominator take the pixel-parallel form
\begin{align}
  \mu_N^{F}
  &\asymp
  \bstar^{\top}\widetilde{\bm\beta}_{\kappa},
  \label{eq:feature-nacc-de}\\
  \mu_D^{F}
  &\asymp
  \left\lVert\widetilde{\bm\beta}_{\kappa}\right\rVert^2
  +\frac{\kappa^2B^{F}_{1,2}+S_{\perp}+\sigma^2}
         {n-\mathrm{df}^{F}_{2}}
   \tr\!\left[
     FF^{\top}(\Sigma_F+\kappa I)^{-2}\Sigma_F
   \right]
  \label{eq:feature-dacc-de}\\
  &=
  \left\lVert\widetilde{\bm\beta}_{\kappa}\right\rVert^2
  +\frac{E_{\rm gen}^{\rm DE}+\sigma^2}{n}
   \tr\!\left[
     FF^{\top}(\Sigma_F+\kappa I)^{-2}\Sigma_F
   \right].
  \label{eq:feature-dacc-egen}
\end{align}
Thus
\begin{equation}
  z_{\rm acc}^{F,\rm DE}
  \asymp \frac{\mu_D^{F}}{\mu_N^{F}}-1,
  \qquad
  R_{\rm acc}^{2,\rm DE}
  \asymp 1-\left(z_{\rm acc}^{F,\rm DE}\right)^2.
\end{equation}
Setting $F=I$ gives
$\Sigma_F=\Sigma$, $\bm w^{\star}=\bstar$, $S_{\perp}=0$, and recovers
Proposition~\ref{prop:pixel-de} exactly.
\end{proposition}

\begin{theorem}[Prediction statistics do not identify pixel control]
\label{thm:summary-separation}
For Gaussian inputs and a deterministic feature map, the joint distribution
of the training features and responses is determined by
\begin{equation}
  \left(C,\bm h,S,\sigma^2\right).
\end{equation}
Consequently, the distribution of the fitted ridge coefficient, the complete
natural generalization curve, and population/DE cross-validated ridge
selection depend only on
\begin{equation}
  \left(C,\bm h,S,\sigma^2,n\right).
\end{equation}
Pixel-gradient accentuation additionally depends on
\begin{equation}
  \left(G,\bm q\right)
  =\left(FF^{\top},F\bstar\right).
\end{equation}
Therefore prediction and cross-validation alone do not identify
$E_{\rm acc}$, $b_{\rm acc}$, or $R_{\rm acc}^{2}$.
\end{theorem}

\begin{itemize}
  \item Present the summary-statistics dependency graph here, immediately
        after Theorem~\ref{thm:summary-separation}.
  \item For non-Gaussian natural inputs, state the weaker exact claim:
        $(C,\bm h,S)$ determine the second-order prediction theory, but need
        not determine the full feature--response distribution.
\end{itemize}

\begin{figure}[t]
  \centering
  \fbox{\parbox{0.93\linewidth}{
    \textbf{Figure 3 placeholder: dependency graph.}
    Primitive quantities $(F,\Sigma,\bstar)$ feed into a prediction branch
    $(C,\bm h,S,\tstar,S_\perp,\kappa_{\rm CV},E_{\rm gen})$ and a control
    branch that additionally requires $(G,\bm q)$. Keep the two branches
    visually separated and left-to-right topologically ordered.
  }}
  \caption{Prediction and pixel control use overlapping but nonidentical
  summary statistics.}
  \label{fig:dependency-graph}
\end{figure}

\section{An exact prediction-preserving family with changing control}
\label{sec:gauge}

\begin{itemize}
  \item Move to whitened input coordinates, assuming $\Sigma$ is positive
        definite on the effective support:
  \begin{equation}
    A=F\Sigma^{1/2},\qquad
    \bm b=\Sigma^{1/2}\bstar,
    \qquad
    \bm c=\Sigma^{-1/2}\bstar.
  \end{equation}
  \item In these coordinates, $C=AA^{\top}$, $\bm h=A\bm b$,
        $\bm q=A\bm c$, while
        $G=A\Sigma^{-1}A^{\top}$. Use support-restricted inverses for singular
        $\Sigma$.
\end{itemize}

\begin{theorem}[Prediction-preserving gauge transformation]
\label{thm:gauge}
Let $Q\in O(d)$ be orthogonal and satisfy $Q\bm b=\bm b$. Define
\begin{equation}
  F_Q=F\Sigma^{1/2}Q\Sigma^{-1/2}
      =AQ\Sigma^{-1/2}.
  \label{eq:gauge-map}
\end{equation}
Then
\begin{equation}
  F_Q\Sigma F_Q^{\top}=F\Sigma F^{\top}=C,
  \qquad
  F_Q\Sigma\bstar=F\Sigma\bstar=\bm h.
\end{equation}
Hence the Gaussian training distribution, ridge estimator in feature
coordinates, natural generalization curve, and cross-validated ridge penalty
are invariant under $Q$. In contrast,
\begin{align}
  \bm q_Q&=F_Q\bstar=AQ\bm c,\\
  G_Q&=F_QF_Q^{\top}=AQ\Sigma^{-1}Q^{\top}A^{\top},
\end{align}
which need not be invariant.
\end{theorem}

\begin{corollary}[Strong gauge: isolate the back-propagation metric]
\label{cor:strong-gauge}
If $Q$ fixes both $\bm b$ and $\bm c$, then $C$, $\bm h$, and $\bm q$ are
all invariant, while $G_Q$ can still change whenever $\Sigma$ is anisotropic.
Along this stronger family, the distribution of
$N_{\rm acc}=\bm q^{\top}\that$ is fixed and all changes in accentuation arise
through the back-propagated gradient energy
$D_{\rm acc}=\that^{\top}G_Q\that$.
\end{corollary}

\begin{itemize}
  \item Explain why the construction sometimes projects away the span of
        both $\bm b$ and $\bm c$. Fixing $\bm b$ alone is sufficient to
        preserve prediction. Fixing $\bm c$ as well preserves the
        accentuation numerator and yields the cleanest causal dissection of
        $G$.
  \item Connect the gauge to anisotropy. If $\Sigma=I$, then
        $G_Q=AA^{\top}$ is invariant and the nontrivial dissociation
        disappears.
\end{itemize}

\subsection{Continuous top-PC-to-tail path used in the experiment}

\begin{itemize}
  \item Use the strong-gauge subspace so every moving component annihilates
        both $\bm b$ and $\bm c$.
  \item Decompose the whitened feature map into a fixed part $P$, a
        well-behaved top-PC component $R$, and a matched tail component $V$:
  \begin{equation}
    A(t)=P+\sqrt{1-t}\,R+\sqrt t\,V,
    \qquad
    F(t)=A(t)\Sigma^{-1/2},
    \qquad 0\le t\le1.
    \label{eq:continuous-path}
  \end{equation}
  \item State the construction constraints:
  \begin{equation}
    R\bm b=V\bm b=0,
    \qquad
    R\bm c=V\bm c=0,
    \qquad
    RR^{\top}=VV^{\top},
  \end{equation}
  together with the required cross-orthogonality to $P$ and between the
  rotating components. These conditions make $A(t)A(t)^{\top}$,
  $A(t)\bm b$, and $A(t)\bm c$ invariant.
  \item Parameterize $t=\sin^2\theta$, so that
        $\sqrt{1-t}=\cos\theta$ and $\sqrt t=\sin\theta$. The path is an
        ordinary continuous rotation in whitened coordinates.
  \item Choose $t=0$ to be the unwhitened top-$p$ PC projector. Choose $V$
        in very low-variance spectral directions, making $G(t)$ grow while
        prediction remains fixed.
\end{itemize}

\begin{proposition}[Protected prediction can approach singular control]
\label{prop:protected-singularity}
Assume $C$ is positive definite on the effective feature support. The smallest
excess natural-prediction loss among feature weights with zero true
accentuation gain is
\begin{equation}
  \epsilon_{N=0}
  =\min_{\bm\theta:\,\bm q^{\top}\bm\theta=0}
     \lVert\bm\theta-\tstar\rVert_C^2
  =\frac{(\bm q^{\top}\tstar)^2}
         {\bm q^{\top}C^{-1}\bm q}.
  \label{eq:cost-zero-gain}
\end{equation}
Thus nearly perfect natural prediction can lie arbitrarily close to the
$R_{\rm acc}^{2}$ singularity when both $S_{\perp}/S$ and
$\epsilon_{N=0}/S$ are small. The susceptibility
$\bm q^{\top}C^{-1}\bm q$ is large when the control functional points into
low-variance feature directions.
\end{proposition}

\begin{figure}[t]
  \centering
  \fbox{\parbox{0.93\linewidth}{
    \textbf{Figure 4 placeholder: main constructive result.}
    Use the Van Hateren $p=500$ top-PC-to-tail rotation. Top row:
    $\tr(FF^{\top})$, invariant $E_{\rm gen}$, and invariant
    $R_{\rm gen}^{2}$. Bottom row: $E_{\rm acc}/S$, $R_{\rm acc}^{2}$, and
    $b_{\rm acc}$ across $t=\sin^2\theta$, with noise levels encoded by color.
    Pair this with selected feature-row spectra at $t=0,10^{-4},10^{-2},1$.
  }}
  \caption{A prediction-equivalent feature rotation moves from a controlled
  top-PC endpoint to catastrophic pixel-space accentuation.}
  \label{fig:gauge-construction}
\end{figure}

\section{Natural-image experiments}
\label{sec:experiments}

\subsection{Measured covariance and disk teacher}

\begin{itemize}
  \item Introduce the Van Hateren and FFHQ image ensembles, patch/image
        dimensions, covariance estimation, retained rank, and circular disk
        teacher.
  \item Plot the image spectrum and teacher power in the population-PC basis.
  \item Report response noise on two aligned axes: response standard deviation
        $\sigma$ and noise-to-signal variance ratio $\sigma^2/S$.
  \item Explain which experiments draw Gaussian samples with the measured
        covariance and which use actual natural-image samples. Do not imply
        exact Gaussian distributional theory for the latter.
\end{itemize}

\subsection{Cross-validated pixel ridge}

\begin{itemize}
  \item Compare fixed ridge and prediction-selected RidgeCV across noise.
  \item Show that CV protects generalization substantially but does not make
        accentuation an optimized objective.
  \item Plot the selected $\alpha$ path beneath vertically aligned panels for
        $E$, $R^2$, and slope so transitions can be related to regularization.
  \item Use the leading DE for typical values and the second-order or
        distributional treatment for nonlinear ratio averages.
\end{itemize}

\subsection{PCA, whitening, and top-PC truncation}

\begin{center}
\begin{tabular}{@{}lllll@{}}
\toprule
Feature map & feature loading & $C$ spectrum & $G$ spectrum & omitted signal\\
\midrule
Full PCA & $1$ & $s_j$ & $1$ & $0$\\
Full whitening & $s_j^{-1/2}$ & $1$ & $s_j^{-1}$ & $0$\\
Top-$K$ PCs & $1$, $j\le K$ & $s_j$ & $1$ & $\sum_{j>K}s_jb_j^2$\\
\bottomrule
\end{tabular}
\end{center}

\begin{itemize}
  \item Full PCA is the direct pixel baseline in rotated coordinates.
  \item Whitening makes regression isotropic in feature space but amplifies
        low-variance directions by $1/s_j$ when gradients are mapped back to
        pixels.
  \item Top-PC truncation pays an irreducible prediction cost but removes
        harmful low-variance control directions. When teacher signal is
        concentrated in top PCs, this trade can improve accentuation greatly
        at modest predictive cost.
\end{itemize}

\begin{figure}[t]
  \centering
  \fbox{\parbox{0.93\linewidth}{
    \textbf{Figure 5 placeholder: representation comparison.}
    Noise sweeps for full PCA, whitening, and top-PC truncation. Overlay
    $R_{\rm gen}^{2}$ and $R_{\rm acc}^{2}$ in one panel, and
    $b_{\rm gen}$ and $b_{\rm acc}$ in another. Include the top-PC cutoff
    sweep if space permits.
  }}
  \caption{Feature conditioning for prediction can worsen pixel-space control;
  spectral truncation can reverse the gap.}
  \label{fig:feature-comparison}
\end{figure}

\section{Finite-size fluctuations of nonlinear control metrics}
\label{sec:distributional-de}

\begin{itemize}
  \item Explain why a ratio of deterministic means can predict typical
        accentuation while failing for an empirical mean near a singular
        denominator.
  \item For this finite-size subsection only, retain the Gaussian approximation
        to the fitted feature coefficient using notation parallel to the
        preceding proposition:
  \begin{align}
    \bm m_F
      &:=\Sigma_F(\Sigma_F+\kappa I)^{-1}\bm w^{\star},\\
    \tau_F
      &:=\frac{\kappa^2B^F_{1,2}+S_{\perp}+\sigma^2}
                   {n-\mathrm{df}^{F}_{2}}
       =\frac{E_{\rm gen}^{\rm DE}+\sigma^2}{n},\\
    V_F
      &:=\tau_F(\Sigma_F+\kappa I)^{-1}
                    \Sigma_F(\Sigma_F+\kappa I)^{-1},\\
    \that
      &\ \dot\sim\ \mathcal N(\bm m_F,V_F).
    \label{eq:gaussian-de}
  \end{align}
\end{itemize}

\begin{proposition}[Gaussian distributional deterministic equivalent]
\label{prop:gaussian-de}
Under Eq.~\eqref{eq:gaussian-de}, the linear--quadratic accentuation moments
are
\begin{align}
  \E N_{\rm acc}&=\bm q^{\top}\bm m_F,
  &\Var(N_{\rm acc})&=\bm q^{\top}V_F\bm q,\\
  \E D_{\rm acc}&=\bm m_F^{\top}G\bm m_F+\tr(GV_F),
  &\Var(D_{\rm acc})
    &=2\tr[(GV_F)^2]+4\bm m_F^{\top}GV_FG\bm m_F,\\
  &&\Cov(N_{\rm acc},D_{\rm acc})
    &=2\bm q^{\top}V_FG\bm m_F.
\end{align}
Sampling $\that$ from this surrogate and applying the exact per-fit metric
definitions yields approximate means, medians, intervals, and tail
probabilities without replacing random ratios by ratios of means.
\end{proposition}

\begin{itemize}
  \item Diagnose denominator concentration using
  \begin{equation}
    \mathrm{CV}_{N}
      =\frac{\sqrt{\Var(N_{\rm acc})}}{|\E N_{\rm acc}|},
    \qquad
    \mathrm{CV}_{D}
      =\frac{\sqrt{\Var(D_{\rm acc})}}{\E D_{\rm acc}}.
  \end{equation}
  \item Distinguish two questions: whether the leading DE is statistically
        adequate, and whether control itself is good. Increasing $p$ can make
        the leading DE more accurate while making accentuation worse.
  \item Keep the full random-design $\Var(R)$ and CV-selection fluctuation
        theory in the supplement. In the main paper, report medians/IQRs as
        well as means/standard errors near ratio singularities.
\end{itemize}

\section{Peer review and cross-model asymmetry}
\label{sec:peer-review}

\begin{itemize}
  \item Let generator and evaluator have pixel weights $\bhat_g$ and
        $\bhat_e$. Along the generator path, define the teacher and evaluator
        gains
  \begin{equation}
    T_g=\bstar^{\top}\bhat_g,
    \qquad
    D_{e\leftarrow g}=\bhat_e^{\top}\bhat_g.
  \end{equation}
  \item The peer-review coefficient of determination is
  \begin{equation}
    R_{e\leftarrow g}^{2}
      =1-\left(1-\frac{D_{e\leftarrow g}}{T_g}\right)^2.
    \label{eq:peer-r2}
  \end{equation}
  \item Explain the asymmetry geometrically: one generator can move along a
        direction that remains legible to a good evaluator, while the reverse
        direction has nearly zero teacher gain.
  \item Treat peer review as a consequence and validation of the
        prediction--control geometry, not as a second independent theoretical
        spine.
  \item Place the full joint-resolvent theory for correlated fits in the
        supplement or future work.
\end{itemize}

\section{Discussion}
\label{sec:discussion}

\begin{itemize}
  \item Main conclusion: predictive equivalence is weaker than control
        equivalence because intervention exposes geometry not identifiable
        from natural prediction alone.
  \item Cross-validation solves the objective it observes. It can select the
        prediction-optimal ridge penalty while leaving substantial or
        catastrophic control error.
  \item Low-variance directions provide a shared mechanism for the pixel and
        feature cases, but feature maps add an irreducible structural mismatch
        between prediction-optimal and control-optimal projections.
  \item Whitening illustrates that a representation can be statistically
        well conditioned for regression while being ill conditioned for
        pixel-space intervention.
  \item Top-PC truncation illustrates a practical design principle: discard
        directions that contribute little target signal but large intervention
        gain.
  \item Relate the linear theory cautiously to nonlinear neural networks:
        local input gradients define an input-space control metric, while
        representation covariance governs prediction on natural data.
  \item State limitations: Gaussian or second-order input model, fixed linear
        features, linear teacher, local gradient paths, and incomplete
        unconditional fluctuation theory for nonlinear ratios and random CV.
\end{itemize}

\section*{Proposed supplementary organization}

\begin{enumerate}
  \item Exact generalization, accentuation, and peer-review metric derivations.
  \item One- and two-resolvent deterministic-equivalent identities.
  \item Per-PC bias and variance decomposition for direct ridge regression.
  \item Derivation of the CV stationarity condition and
        Eq.~\eqref{eq:pixel-z-cv}.
  \item Arbitrary fixed-feature-map DE and aligned spectral specialization.
  \item Proofs of Theorems~\ref{thm:summary-separation} and~\ref{thm:gauge},
        including the continuous rotation construction.
  \item Delta-method corrections, conditional $\Var(R)$, missing
        random-design terms, and CV-selection variability.
  \item Gaussian distributional DE and effective-rank diagnostics.
  \item Additional dimensions, random null-space rotations, FFHQ replication,
        and sensitivity analyses.
\end{enumerate}

\section*{Open author decisions before full drafting}

\begin{itemize}
  \item Choose whether the primary empirical opening is the original neural
        control result or the natural-image disk-teacher demonstration.
  \item Choose one primary natural-image ensemble for the main text; place the
        second ensemble in the supplement as replication.
  \item Decide whether peer review receives a main figure panel or only a
        concise corollary plus supplementary validation.
  \item Decide whether the main paper reports both $E_{\rm acc}/S$ and
        $R_{\rm acc}^{2}$ throughout. Keeping both is theoretically useful
        because they expose opposite ratio singularities.
  \item Decide how much of the Gaussian distributional DE belongs in the main
        Results versus Methods/Supplement.
  \item Select a target venue and length before converting the bullets to
        complete prose.
\end{itemize}

\end{document}
